% Chapter 1, section: orientation
\section{Orientation in 3-D}\label{sec:rot:orientation}

\subsection{Orientation is a relation between two frames}\label{sec:rot:orientation-relation-between}
A single direction vector (say, where the nose points) is \emph{not} an orientation: you can still spin the body about that direction. To pin the body down you need its whole frame $\{b_x,b_y,b_z\}$, and that frame only means something relative to another frame $\{a_x,a_y,a_z\}$. Orientation is therefore a \emph{relation between two frames}, and the natural object to store is: ``the body axes, written in reference coordinates''.

\subsection{The rotation matrix from first principles}\label{sec:rot:rotation-matrix-first}
\begin{figure}[h]
\centering\includegraphics[width=0.95\linewidth]{frames.pdf}
\caption{(a) Reference frame $\{A\}$ (grey) and body frame $\{B\}$ (colour) sharing an origin. (b) The first column of $R_{AB}$ is the arrow $b_x$ decomposed along $a_x,a_y,a_z$; each component is a projection $r_{i1}=a_i\cdot b_x$.}
\label{fig:rot:frames}
\end{figure}
Write each body axis in $A$-coordinates and stack them as columns:
\begin{equation}
R_{AB}=\Big[\ \pre{A}{b_x}\ \ \pre{A}{b_y}\ \ \pre{A}{b_z}\ \Big]
=\bmat{a_x\!\cdot b_x & a_x\!\cdot b_y & a_x\!\cdot b_z\\ a_y\!\cdot b_x & a_y\!\cdot b_y & a_y\!\cdot b_z\\ a_z\!\cdot b_x & a_z\!\cdot b_y & a_z\!\cdot b_z},\qquad r_{ij}=a_i\cdot b_j=\cos\angle(a_i,b_j).
\label{eq:rot:dcm}
\end{equation}
This is why rotation matrices are also called \emph{direction-cosine matrices}. Why does it convert coordinates? Take any vector with body coordinates $\pre{B}{p}=(p_1,p_2,p_3)$, meaning $p=p_1b_x+p_2b_y+p_3b_z$. Writing each $b_j$ in $A$-coordinates and summing:
\[
\pre{A}{p}=p_1\,\pre{A}{b_x}+p_2\,\pre{A}{b_y}+p_3\,\pre{A}{b_z}=R_{AB}\,\pre{B}{p}.
\]
That is literally the definition of matrix--vector multiplication (a weighted sum of columns). Nothing else is needed.

\begin{keyidea}
\emph{Column $j$ of $R_{AB}$ is axis $j$ of $\{B\}$ seen from $\{A\}$; row $i$ of $R_{AB}$ is axis $i$ of $\{A\}$ seen from $\{B\}$.} If you can say which frame's axes are in the columns, you will never confuse $R_{AB}$ with $R_{BA}$.
\end{keyidea}

\subsection{Elementary rotations}\label{sec:rot:elementary-rotations}
\begin{deriv}{$R_z(\theta)$, $R_x(\theta)$, $R_y(\theta)$}
Rotate by $\theta$ about $z$ (counter-clockwise looking down $+z$). From \figref{fig:rot:activepassive}(a), $e_1\mapsto(\cos\theta,\sin\theta,0)$, $e_2\mapsto(-\sin\theta,\cos\theta,0)$, and $e_3$ is on the axis so it stays. Those images are the columns:
\[
R_z(\theta)=\bmat{c_\theta&-s_\theta&0\\ s_\theta&c_\theta&0\\0&0&1},\quad
R_x(\theta)=\bmat{1&0&0\\0&c_\theta&-s_\theta\\0&s_\theta&c_\theta},\quad
R_y(\theta)=\bmat{c_\theta&0&s_\theta\\0&1&0\\-s_\theta&0&c_\theta}.
\]
$R_x$ follows from $R_z$ by the cyclic relabelling $x\to y\to z\to x$ (which preserves handedness). For $R_y$ the cyclic pair is $(z,x)$, not $(x,z)$, because $z\times x=+y$. Writing the block in the order $(x,z)$ therefore moves the minus sign to the lower-left. This is not a typo to memorise; it is the right-hand rule.
\end{deriv}

\subsection{Why $R^{-1}=R\T$, why $\det R=+1$: the group $\SO$}\label{sec:rot:group}
\begin{deriv}{the defining properties of a rotation}
\textbf{(1) Orthogonality.} A physical rotation preserves lengths and angles, hence dot products: $(Rx)\cdot(Ry)=x\cdot y$ for all $x,y$. Then $x\T(R\T R)y=x\T y$ for all $x,y$, which forces $R\T R=I$. Geometrically: the columns are the body axes, which are orthonormal.

\textbf{(2) Inverse is the transpose.} From $R\T R=I$. Read via \eqref{eq:rot:dcm}: $(R_{AB})\T$ has entries $b_i\cdot a_j$, i.e.\ it is ``$A$'s axes written in $B$'', which is $R_{BA}$. So $R_{BA}=R_{AB}^{-1}=R_{AB}\T$: inverting a rotation costs nothing.

\textbf{(3) Determinant $+1$.} $\det(R\T R)=(\det R)^2=1$. A physical rotation can be reached continuously from the identity (just turn the body slowly). $\det$ is continuous and can only take the values $\pm1$, so it cannot jump: it stays at $\det I=1$. Equivalently $\det R=b_x\cdot(b_y\times b_z)=+1$ says the body frame is right-handed.
\end{deriv}
Hence
\[
\boxed{\ \SO=\{R\in\RR^{3\times3}:\ R\T R=I,\ \det R=1\}\ }
\]
is a \textbf{group}: closed under products ($(R_1R_2)\T R_1R_2=I$, $\det=1\cdot1$), has an identity, has inverses ($R\T$), is associative, but is \emph{not} commutative (\S\ref{sec:rot:3-d-rotations}). It is \emph{not a vector space}: $R+R=2R$ has determinant 8 and stretches vectors, and $\frac12(R_1+R_2)$ is generally not a rotation. Averaging, interpolating or adding small corrections therefore needs the tools of \S\ref{sec:rot:exponential-coordinates}.

\subsection{Composition, and the order of multiplication}\label{sec:rot:composition-order-multiplication}
Chaining frames is cancellation of inner indices: $\pre{A}{p}=R_{AB}\pre{B}{p}=R_{AB}R_{BC}\pre{C}{p}\Rightarrow R_{AC}=R_{AB}R_{BC}$. When you instead \emph{apply an extra rotation} $\Delta$ to a body with current orientation $R$, it matters which frame $\Delta$ is described in:
\begin{deriv}{pre-multiply = fixed axes, post-multiply = body axes}
\emph{World axes.} Rotating the body by $\Delta$ about world-fixed axes rotates each column (each body axis): $R_{\text{new}}=\Delta R$.
\emph{Body axes.} ``Rotate by $\Delta$ about the body's own axes'' means: the axis vector $k$ given in body coordinates is $Rk$ in world coordinates. By \eqref{eq:rot:conj}, rotation about $Rk$ is the conjugate $R\,\Delta\,R\T$. Applying that on the left: $R_{\text{new}}=(R\Delta R\T)R=R\Delta$.
\end{deriv}
\textbf{Rule:} pre-multiplication rotates about the fixed (world) axes; post-multiplication rotates about the current (body) axes. This single fact explains intrinsic vs.\ extrinsic Euler angles in \S\ref{sec:rot:euler-angles} and left vs.\ right perturbations in \S\ref{sec:rot:exponential-coordinates}.

\subsection{Relative orientation}\label{sec:rot:relative-orientation}
Given two orientations in a common frame, the rotation between them follows from index cancellation:
\[
R_{BC}=R_{BA}R_{AC}=R_{AB}\T R_{AC}.
\]
The ``size'' of this relative rotation (the angle needed to align the frames) is computed in \S\ref{sec:rot:exponential-coordinates} by $\theta=\arccos\big((\tr R_{BC}-1)/2\big)$.

\subsection{Degrees of freedom: why 3, not 9}\label{sec:rot:degrees-freedom-3}
Nine entries, but $R\T R=I$ is symmetric, so it gives $6$ independent scalar equations (3 unit lengths, 3 perpendicularities). $9-6=3$. Constructive count: choose the direction of $b_x$ (a point on the unit sphere: 2 DOF), then $b_y$ in the circle perpendicular to it (1 DOF); $b_z=b_x\times b_y$ is forced. Any representation with more than 3 numbers carries constraints (quaternion: 1, matrix: 6). Any representation with exactly 3 numbers must have a singularity somewhere (a topological fact about $\SO$ that \S\ref{sec:rot:euler-angles} makes concrete).

\subsection{Rotation vs.\ change of coordinates}\label{sec:rot:rotation-change-coordinates}
The same matrix answers two questions. Let $R=R_z(90^\circ)$ and $p=(1,0,0)$. Active: rotating the vector $p$ by $90^\circ$ about $z$ gives $(0,1,0)$. Passive: if frame $\{B\}$ is $\{A\}$ turned $90^\circ$ about $z$, the point whose $B$-coordinates are $(1,0,0)$ has $A$-coordinates $(0,1,0)$. Same arithmetic, two stories. Write the frames as subscripts and you never need to remember which story a library intended: check that inner indices match.

\begin{example}{camera mounted on a robot (frame conventions)}
A mobile robot uses the ROS body convention ($x$ forward, $y$ left, $z$ up). Its camera uses the optical convention ($z_C$ forward along the optical axis, $x_C$ right in the image, $y_C$ down). The camera looks straight ahead. Find $R_{BC}$, then convert a detection $\pre{C}{p}=(0.2,\,-0.1,\,3.0)$\,m into the body frame.

\emph{Columns = camera axes in body coordinates:} $x_C$ (right) $=-y_B\to(0,-1,0)$; $y_C$ (down) $=-z_B\to(0,0,-1)$; $z_C$ (forward) $=x_B\to(1,0,0)$:
\[
R_{BC}=\bmat{0&0&1\\-1&0&0\\0&-1&0},\qquad \det R_{BC}=+1\ \checkmark,\qquad
\pre{B}{p}=R_{BC}\,\pre{C}{p}=(3.0,\,-0.2,\,0.1).
\]
Sanity check: 3\,m ahead, 0.2\,m to the right ($y_B<0$), 0.1\,m up (the detection was above the image centre, $y_C<0$). Writing this matrix by inspection of columns is far less error-prone than composing ``$-90^\circ$ about $z$ then $-90^\circ$ about $x$'', and it is the first thing to check when a perception pipeline produces mirrored or swapped detections.
\end{example}

\begin{example}{relative orientation of two IMUs (knee angle), part 1}
Two IMUs are strapped to a thigh (1) and a shank (2). Each fusion filter reports its orientation relative to a common world frame. The physiotherapist wants the knee angle, which must not depend on which way the person is facing. With $R_{W1}=R_z(40^\circ)R_y(10^\circ)$ and $R_{W2}=R_{W1}R_y(55^\circ)R_x(3^\circ)$ (numerically, rounded):
\[
R_{W1}=\bmat{0.7544&-0.6428&0.1330\\0.6330&0.7660&0.1116\\-0.1736&0&0.9848},\quad
R_{12}=R_{W1}\T R_{W2}=\bmat{0.5736&0.0429&0.8180\\0&0.9986&-0.0523\\-0.8192&0.0300&0.5728}.
\]
The heading $40^\circ$ has cancelled exactly, as it must: $R_{12}$ is a relation between the two segments only. We will extract the knee angle and hinge axis from $R_{12}$ in \S\ref{sec:rot:log} (Log map). Note the common mistake: subtracting Euler angles reported by the two filters. It works only by luck when the rotations share an axis.
\end{example}

\begin{pitfall}[s in practice]
\begin{itemize}
\item \textbf{Transpose confusion.} If an object appears to rotate the ``wrong way'' or a transform chain only works at zero angle, you used $R_{BA}$ where $R_{AB}$ was needed. Check: the columns should be the axes of the frame named \emph{second}.
\item \textbf{Drift off $\SO$.} After many floating-point products, $R\T R\neq I$ slowly. Repair by projecting to the nearest rotation: with SVD $M=U\Sigma V\T$, use $U\,\mathrm{diag}(1,1,\det(UV\T))\,V\T$. Better: integrate in a way that never leaves $\SO$ (\S\ref{sec:rot:angvel}).
\item \textbf{Mirrored mounts.} A matrix built from hand-measured axes with $\det=-1$ means one axis is flipped. Always check the determinant.
\end{itemize}
\end{pitfall}
